\chapter{Professional Communication: Oral}

\section{Design Review Presentations}
Structure every design review around a \textbf{narrative arc}: problem $\to$ approach $\to$ results $\to$ next steps. The audience should feel like they are following a logical story, not watching a random sequence of slides. Time your presentation and rehearse (going over time signals poor preparation; going significantly under time signals insufficient depth).

Coordinate team transitions: each speaker introduces the next (``Now Sarah will walk you through our concept selection process''), and all team members use a consistent slide template and visual style. A presentation where each section looks like it was made by a different team in a different course undermines credibility.

\section{Slide Design Best Practices}
\begin{itemize}[nosep,leftmargin=1.5em]
\item \textbf{One message per slide}. If you need more than 30 seconds to explain a slide, it has too much content.
\item \textbf{Use visuals, not text walls}. Show the system architecture diagram, the decision matrix, the test data plot. Do not put the full text of your requirements on a slide---show a summary table.
\item \textbf{Data visualization}: choose the right chart type. Bar charts for comparisons. Line charts for trends. Scatter plots for correlations. Tables for exact values. Label all axes with units.
\item \textbf{Font size}: minimum 18pt for body text, 24pt for titles. If you have to say ``I know this is hard to read,'' the slide has failed.
\end{itemize}

\section{Handling Q\&A}
Q\&A is not a threat; it is the most valuable part of the review. Tough questions from your PA, client, or TA improve the design. Listen to the entire question before responding. If you do not understand, ask for clarification (``Could you clarify what you mean by `thermal margin'?''). If you do not know the answer, say so honestly and commit to following up: ``That is an excellent question. I do not have the data at hand, but I will have an answer by our next meeting.'' Never make up an answer.

\section{The Elevator Pitch}
Every team member should be able to explain the project in \textbf{60 seconds} to a non-technical audience. Structure: what problem you are solving (one sentence), for whom (one sentence), how your approach works (one sentence), and why it matters (one sentence). Practice until it is natural. You will use this at Career Day, at the Showcase, and in every job interview for the next decade.

\section{Sprint Review Presentations}
Sprint reviews (Chapter~11) are less formal than design reviews but still require preparation. Show \emph{working functionality}: a sensor reading on the screen, a motor spinning, a simulation result that matches the benchmark. ``We are still working on it'' is not a sprint review; it is an excuse.

\section{Video Demonstrations}
For subsystem demos (SDI) and system demos (SDII), you will record narrated video demonstrations. Narrate what you are doing, what the expected result is, and what the actual result is. A silent video of a blinking LED tells the viewer nothing. Verify your video link is accessible by opening it in an incognito/private browser window. Unviewable links may receive a zero.

\textbf{Good demonstration video practices}: (1)~plan the demo script before recording; what will you show, in what order, and what will you say at each step, (2)~start with a 10-second introduction stating your name, team, and what you are demonstrating, (3)~show the test setup with a wide shot, then zoom in on the relevant component, (4)~narrate the expected behavior before triggering it (``When I apply 5\,V to the input, we expect to see approximately 3.3\,V at the output''), (5)~show the actual result clearly (point to the multimeter reading, show the oscilloscope trace), (6)~state whether the result matches expectation and what it means.

\textbf{Recording tips}: use a phone or webcam with a tripod (not handheld; shaky video is unwatchable), ensure adequate lighting, minimize background noise, and keep videos under 3 minutes per subsystem. Longer is not better; concise and clear is.

\section{Preparing for Design Reviews: The Rehearsal}
Every design review presentation should be rehearsed at least twice:

\textbf{First rehearsal} (internal, 4--5 days before): run through the full presentation with the team. Time each section. Identify: slides that are too dense (split them), sections that run long (cut or condense), transitions that are awkward (add bridging sentences), and questions you cannot answer (research before the formal review).

\textbf{Second rehearsal / Mock Review} (with PA or another team, 2--3 days before): present to fresh eyes. Note what questions they ask; these are the same questions your formal reviewers will ask. Practice answering concisely: a 30-second answer is almost always better than a 3-minute answer for a review question.

\textbf{Day-of preparation}: arrive 15 minutes early. Test the projector, screen, and any live demonstration equipment. Have backup slides (PDF on a USB drive) in case the primary laptop fails. Assign one team member to advance slides so the presenter can focus on speaking.

\section{Managing Presentation Nervousness}
Nervousness before a presentation is universal and physiologically normal (elevated heart rate, sweaty palms, dry mouth). It is not a sign of inadequacy; it is a sign that you care about the outcome. Strategies:

\textbf{Preparation is the best antidote}: the single most effective way to reduce nervousness is thorough preparation. If you know your material deeply and have rehearsed twice, your confidence is grounded in competence. If you are reading from slides for the first time, your nervousness is grounded in justified uncertainty.

\textbf{Structured breathing}: before stepping up to present, take 4--5 slow, deep breaths (4 seconds in, 4 seconds hold, 6 seconds out). This activates the parasympathetic nervous system and reduces the fight-or-flight response.

\textbf{Focus outward, not inward}: instead of worrying about how you look or sound, focus on the audience. Are they understanding the content? Are they engaged? Do they have questions? Outward focus redirects attention from self-consciousness to communication.

\textbf{Accept imperfection}: you will stumble on a word, lose your place, or forget a detail. This is normal. Pause, take a breath, and continue. The audience cares about the engineering content, not about flawless delivery.

\section{Coordinating Team Presentations}
A design review presentation with 4--5 speakers requires coordination:

\textbf{Transitions}: each speaker introduces the next (``Now Sarah will walk through our concept selection process''). Never end a section with silence and a confused handoff.

\textbf{Visual consistency}: all slides should use the same template, font, color scheme, and labeling conventions. A presentation where each section looks like it was designed by a different person undermines the impression of a cohesive team.

\textbf{Q\&A protocol}: agree in advance who answers questions about which topics. If a question is directed to the wrong person, the Scrum Master redirects: ``That's a great question about the power subsystem; let me pass it to Jordan, who led that design.'' Never have two people start answering the same question simultaneously.

\textbf{Time management}: assign a timekeeper (not the current speaker). At the 2-minute warning, the timekeeper signals. At time, the speaker wraps up the current slide and transitions. Going over time cuts into Q\&A, which is the most valuable part of the review.

\section{Synthesizing Feedback}
After every design review, consolidate all feedback (from PA, client, TAs, and peers) into an \textbf{action item list}. For each item: what was the feedback, what action is required, who is responsible, and when will it be completed. Report back on closed action items at the next sprint review. Feedback that is received but not acted upon is wasted effort by everyone involved.


\section{Speaking to Different Audiences}
The same project must be communicated differently to different audiences:

\textbf{Technical audience} (PA, TAs, fellow engineers): emphasize methodology, analysis rigor, quantitative results, and engineering trade-offs. Use technical terminology. Show equations, data plots, and detailed block diagrams. Expect probing questions about assumptions and limitations.

\textbf{Client audience} (especially non-technical clients): emphasize the problem being solved, the user experience, the key performance outcomes, and the value delivered. Minimize jargon. Show photos, demonstrations, and high-level block diagrams. Translate technical metrics into client-meaningful terms: ``the sensor accuracy is $\pm$0.5\,\degC{}'' becomes ``the system will detect temperature changes as small as one degree, which is sufficient to prevent crop damage in your greenhouse.''

\textbf{General audience} (Showcase visitors, prospective students, administrators): emphasize the ``so what''; why does this project matter? What real-world problem does it address? What makes the solution interesting or innovative? Keep it accessible. A visitor should understand and appreciate your project in 2 minutes.

\textbf{Adapting in real time}: at design reviews and the Showcase, you will encounter all three audiences in rapid succession. Develop the ability to shift your communication register based on the listener's questions. If they ask ``What's your control bandwidth?'' they are technical; go deep. If they ask ``What does it do?'' they are general; go broad.

\section{Using Visual Aids Effectively}
Visual aids (slides, posters, demos, videos) should support your message, not be your message. If you are reading your slides verbatim, the slides have replaced you as the presenter. Best practices:

\textbf{The 6-word test}: each slide should have a clear takeaway that can be stated in 6 words or fewer. If you cannot summarize the slide's message concisely, the slide is doing too much.

\textbf{Build animations}: for complex diagrams or data, reveal information progressively rather than showing everything at once. A block diagram that builds from left to right mirrors your verbal explanation.

\textbf{Live demonstrations}: nothing is more compelling than showing the system working in real time. But live demos can fail. Always have a backup video of the demonstration in case the hardware misbehaves. Test the demo in the presentation room before the review.

\textbf{Pointing}: when presenting data or diagrams, physically point to the relevant area of the slide. This directs the audience's attention and prevents them from trying to parse the entire figure while you are talking about one specific element.



\section{The Showcase: Speaking to a General Audience}
The Design Showcase is your public exhibition. Visitors include: other students, faculty, industry representatives, prospective students and their parents, university administrators, and media. Your communication must be accessible to all of them.

\textbf{The 60-second version}: have a memorized, natural-sounding overview that any visitor can understand. Structure: ``We built [what] for [whom] because [why]. It works by [how, one sentence]. The key result is [what you achieved].'' Practice until it flows naturally without sounding rehearsed.

\textbf{The 3-minute version}: for visitors who are interested and want depth. Walk them through the poster: problem (top left), approach (center), results (right), and impact (bottom). Point to the relevant section as you speak. Show the prototype if possible.

\textbf{The technical deep-dive}: for PAs, faculty, or industry engineers who ask detailed questions. Be prepared to discuss: your analysis methodology, specific design decisions and trade-offs, test results with quantitative detail, and lessons learned. Not every team member needs to be prepared for deep-dive questions on every topic; but every subsystem should have at least one team member who can go deep.

\textbf{Handling ``dumb'' questions}: no question is dumb at the Showcase. A visitor who asks ``What's an Arduino?'' is giving you the opportunity to explain your project to a non-engineer; a skill you will use throughout your career. Answer respectfully and use it as a chance to practice accessible communication.

\section{Panel Presentations and Judged Events}
Some capstone projects involve panel presentations (competition judging, client review boards, or external advisory panels). Panel communication differs from standard presentations:

\textbf{Time discipline is critical}: panels see many presentations in one session. Going over time is disrespectful to the panel and to the teams waiting after you. Practice to fit within the allocated time with 1--2 minutes of margin.

\textbf{Address the entire panel}: make eye contact with different panel members throughout the presentation, not just the one who asked the last question. Panels evaluate how well you communicate, not just what you communicate.

\textbf{Anticipate panel-specific questions}: if the panel includes a safety officer, expect safety questions. If it includes a budget manager, expect cost questions. Research the panel composition (if available) and prepare accordingly.

\textbf{Defer gracefully}: if a panel member asks a question outside your expertise, say ``That is outside my area; let me pass it to [teammate] who led that subsystem.'' This is professional, not weak.



\section{Advanced Q\&A Strategies}
Q\&A is where the most valuable learning happens at design reviews. Prepare for it:

\textbf{Anticipate questions}: before every review, brainstorm 10--15 questions that reviewers are likely to ask. For each, prepare a concise answer (30--60 seconds). Common question categories: ``Why did you choose X over Y?'' (design decisions), ``What happens if X fails?'' (risk/failure modes), ``How do you know X works?'' (verification evidence), ``What is the accuracy/uncertainty of X?'' (measurement confidence), and ``What would you do differently?'' (self-awareness).

\textbf{The structured answer}: use the PREP framework for answering technical questions: \textbf{P}oint (state your answer in one sentence), \textbf{R}eason (explain why), \textbf{E}vidence (cite data, analysis, or test results), \textbf{P}oint (restate the answer). This structure prevents rambling and ensures the reviewer gets a complete answer.

\textbf{When you do not know}: ``That is a great question. I do not have the data to answer it right now, but I will investigate and have an answer by [specific date].'' Record the question on the feedback log. Follow through. A reviewer who asks a question and receives a follow-up email with the answer within a week remembers a professional, competent team.

\textbf{When the question challenges your approach}: do not become defensive. Listen completely. Acknowledge the concern: ``You raise a valid point about thermal expansion.'' Then either (a)~explain the analysis you have already done that addresses the concern, or (b)~acknowledge that you have not considered it and commit to investigating: ``We have not analyzed the thermal expansion effect. We will add it to our Sprint~12 analysis plan and report back at FDR.'' Both responses are professional; defensiveness is not.

\textbf{When a teammate gives a wrong answer}: do not contradict them publicly. After the review, discuss privately. If the incorrect information could affect a decision, send a correction in the post-review Next Steps Email: ``Upon further review, we want to clarify that the motor current draw is 2.4\,A, not 1.4\,A as stated during Q\&A. This does not change our power supply selection, which is rated for 3\,A.''

\section{Remote and Hybrid Presentations}
When presenting remotely (Zoom, Teams) or to a hybrid audience:

\textbf{Technology check}: test your setup 30 minutes before the presentation. Verify: camera, microphone (use a headset, not laptop speakers), screen sharing, and internet connection. Have a backup plan: if screenshare fails, email the PDF to the meeting organizer and they share their screen.

\textbf{Engagement}: look at the camera, not the screen (this creates the impression of eye contact). Speak slightly slower and more clearly than in person; audio compression and latency make fast speech harder to follow. Pause after key points to allow questions.

\textbf{Slide design for remote}: increase font sizes by 20\% compared to in-person slides. Remote audiences view on smaller screens (laptops, tablets) and cannot lean forward to read fine print. Use high-contrast colors; low-contrast slides that look fine on a projector may be unreadable on a compressed video stream.

\textbf{Managing Q\&A remotely}: designate one team member to monitor the chat for questions while another presents. At the end of each section, the chat monitor relays questions: ``We have a question from Dr.~Smith in the chat about our sensor placement.'' This ensures remote reviewers are heard.



\section{Turning Nervousness into Energy}
The physiological symptoms of nervousness (elevated heart rate, adrenaline, heightened alertness) are identical to the symptoms of excitement. The difference is interpretation. Research shows that reframing anxiety as excitement (``I am excited to present our work'') improves performance more than trying to calm down (``I need to relax'')~[41].

\textbf{Practical techniques}:
\begin{itemize}[nosep,leftmargin=1.5em]
\item \textbf{Power posing}: stand in an expansive, confident posture for 2 minutes before presenting. Research on this is debated, but many presenters report feeling more confident.
\item \textbf{Warm-up}: say your first 30 seconds aloud in the hallway before entering the room. The hardest part of any presentation is the first sentence; once you are moving, momentum carries you.
\item \textbf{Memorize the opening}: know your first 2--3 sentences by heart. This eliminates the ``blank mind'' moment that novice presenters fear. Once past the opening, your preparation and knowledge of the material take over.
\item \textbf{Focus on the message, not on yourself}: you are not performing; you are communicating engineering work that matters. The audience wants to understand your project, not evaluate your presentation style. Shift your mental focus from ``how do I look?'' to ``did the audience understand the key point?''
\end{itemize}

\section{The Team Presentation Rehearsal Protocol}
For design review presentations with 4--5 speakers:

\textbf{Rehearsal 1} (5 days before, 90 min): full run-through with timing. No interruptions. Record video if possible. After: note time per section, identify over-long sections, flag missing content, and list slides that need redesign.

\textbf{Rehearsal 2} (3 days before, 60 min): revised slides, focused on: transitions between speakers (smooth, with introductions), Q\&A practice (team members ask each other tough questions), and time management. If still over time, cut content from the least essential sections.

\textbf{Rehearsal 3} (day before, 30 min, optional): final dry run in the actual presentation room if possible. Test technology (projector, demo equipment). Walk through the first 2 minutes of each section to verify timing.

The investment of 3 hours in rehearsal consistently produces: a 10--15\% improvement in reviewer-assessed presentation quality, better time management, smoother transitions, and more confident Q\&A responses. This is the highest-ROI activity in design review preparation.



\section{Data Visualization for Design Review Slides}
Engineering presentations live and die on the quality of data visualization:

\textbf{The ``so what'' title}: every data slide should have a title that states the conclusion, not just the topic. Bad: ``Temperature Test Results.'' Good: ``System Maintains $\pm$0.7\,\degC{} Accuracy, Meeting the $\pm$1\,\degC{} Requirement.'' The title tells the reviewer what to conclude from the data before they even look at the plot.

\textbf{Comparison plots}: when presenting test results against requirements, show both on the same plot. A horizontal dashed line at the requirement threshold, with data points above or below, instantly communicates pass/fail. Color-code: green for passing, red for failing.

\textbf{Before/after comparisons}: when presenting design improvements (``we reduced noise by implementing star grounding''), show the before and after data side-by-side on the same scale. The visual contrast communicates the improvement more effectively than a percentage.

\textbf{Simplify for the audience}: a detailed MATLAB plot with 6 data series, 3 axis labels, and a legend in 8-point font may be scientifically complete but is unreadable on a projected slide. Simplify: show only the most relevant data series, increase font sizes, use distinct colors (not shades of the same color), and remove gridlines that do not aid interpretation.

\textbf{Animation for complex data}: for time-series data, sensor comparison, or process flow, consider animating the data (reveal data points sequentially) rather than showing everything at once. This guides the audience's attention through the narrative: ``Here is the setpoint change at $t=30$\,s... now watch how the system responds... settling time is 8 seconds.''

\textbf{Tables vs. plots}: use tables for exact values that the audience needs to reference (BOM summary, requirement compliance matrix). Use plots for trends, comparisons, and patterns. Never put a 50-row table on a slide and say ``the key number is in the third column, eighth row.''



\section{Common Presentation Mistakes and Fixes}
These patterns appear at nearly every design review cycle:

\textbf{The text wall}: slides with 8+ bullet points of full sentences. Fix: convert text to visuals. A requirements slide should show a summary table, not paragraphs. A design slide should show a block diagram, not a description. The presenter provides the narrative; the slides provide the visual evidence.

\textbf{The speed run}: the team practices once, discovers they are 5 minutes over, and decides to ``just talk faster'' instead of cutting content. Fix: remove slides. Identify the 3--5 slides that contribute least to the narrative and eliminate them. A concise 20-minute presentation is always better than a rushed 25-minute presentation.

\textbf{The disconnected Q\&A}: the team gives a polished presentation, then falls apart during Q\&A because members cannot answer questions about sections they did not write. Fix: every team member reviews the entire presentation and can speak to the key points of every section, even if they did not author it.

\textbf{The invisible transition}: the team changes speakers without any verbal bridge, leaving the audience confused about who is speaking and what the new topic is. Fix: every transition includes: ``Thank you, [previous speaker]. Now [next speaker] will walk us through [topic].'' This takes 5 seconds and dramatically improves flow.

\textbf{The apologetic opener}: ``Sorry this slide is hard to read'' or ``We didn't have time to finish this analysis.'' Fix: if the slide is hard to read, redesign it before the review. If the analysis is incomplete, present what you have and state the plan to complete it; without apologizing. Confidence in what you \emph{have} done is more professional than apology for what you have not.


\section{Practice Questions}
\begin{enumerate}[leftmargin=1.5em]
\item Write a 60-second elevator pitch for your capstone project (or a hypothetical one).
\item A reviewer at PDR says ``I don't think your sensor has adequate range for the operating conditions.'' You are not sure whether their concern is valid. How do you respond in the moment, and what follow-up do you take?
\item Your team's CDR presentation runs 8 minutes over the allotted time. Diagnose three likely causes and propose fixes for the FDR presentation.
\end{enumerate}


